\scanpage{42}
\deleted{\diagram{\begin{tikzcd}[ampersand replacement=\&, column sep=large,row sep=large]
f^{-1}(U)\arrow[r,hook,"i"]\arrow[d,"f"']\&X\arrow[d,"f"]\\
U\arrow[r,"j"']\arrow[ur,"g"]\arrow[ur,bend left=12,"\deleted{\unsure{g}}"',no head]\&Y
\end{tikzcd}}}
\deleted{\[gf=i\qquad fg=j\]}
\deleted{$f$ morphisme birationnel \inserted{propre} de préschémas intègres $X,Y$

$U$ ouvert de $Y$ sur lequel $f^{-1}$ est défini

[Si $Y$ est régulier en codim 1, alors $U$ est le complémentaire d'un ens de codim $\geq2$], \deleted{\unsure{i.e. aussi le ft}} \unsure{i.e.} \unsure{aussi le} \unsure{ft|$\forall$} \unsure{que} un \unsure{voisinage} $V$ \unsure{tel} que $f$ induise un isom $f^{-1}(U)\mathbin{-}V$.

Soit $\xi\in\operatorname{Pic}(X)$ défini par un \deleted{\unsure{cycle}} \uline{diviseur} $D$ (\unsure{X connexe}) considérons \deleted{\unsure{un}} \deleted{le morphisme} $g^{-1}(\xi)\in\operatorname{Pic}(U)$, il correspond à \deleted{\unsure{un}} cycle divisoriel}
\margin{41}
